\section{Future design goals}\label{sec:future}
\stagebox{\statusplanned}{none yet}

These rules are deferred. They are recorded here so later work starts from the
same intent.

\begin{enumerate}[leftmargin=*]
\item \textbf{Doubleness.} Stamens turn into petals and fill the centre with
  small, crumpled petals, as in fully double peonies.
\item \textbf{Persistent lengthwise growth.} The growth automaton keeps a
  patch's growth running along the petal's length over several stages, so
  long pleats form instead of only local crumples. Partly built: growth noise
  now persists between stages and can favour length or width
  (\cref{sec:growth}), but nothing yet steers a patch along the petal.
\item \textbf{Spiral or ring placement.} Primordia placed at the golden angle
  ($\approx 137.5^\circ$) in a spiral, or in rings (whorls), chosen as a
  species rule.
\item \textbf{Stamens.} Whether the flower has stamens, and how many, becomes a
  species trait with Gaussian noise.
\end{enumerate}
In the three-tier scheme (\cref{sec:tiers}) stamens are species-wide, with
noise drawn per plant or per flower. Bud shape and core size, once listed here,
are now built (\cref{sec:shells}).
